\documentclass{article}
\usepackage[utf8]{inputenc}
\usepackage{amsmath,amssymb}
\usepackage[most]{tcolorbox}
\usepackage{geometry}
\geometry{margin=2.5cm}

\newcommand{\step}[1]{\par\noindent\hangindent=3.5em\hangafter=1%
  \makebox[3.5em][l]{\textbf{#1.}}}
\newcommand{\steptag}[1]{\hfill{\small\textsf{[#1]}}}

\newtcolorbox{proofbox}[1]{
  colback=gray!5, colframe=gray!60,
  fonttitle=\bfseries, title={#1},
  breakable, enhanced,
  left=4pt, right=4pt, top=4pt, bottom=4pt,
  before skip=10pt, after skip=10pt
}

\begin{document}

\begin{proofbox}{There are universal C\_s2 >= 1 and e\_s2 > 0, with the e\_ca...}
\step{1} There are universal C\_s2 >= 1 and e\_s2 > 0, with the e\_ca threshold and universal C\_ca coefficient of lem-compcb-corner-algebra absorbed into e\_s2 and C\_s2, such that, if A is a finite-dimensional extended epsilon\_A-C*-algebra with epsilon\_A <= t, a supplied MAIN partition state comes from a non-unital extended t-inclusion w:C\textasciicircum{}m->A with one-dimensional images P\_j, has nonempty U, j notin U, dim S\textasciicircum{}A\_\{P\_k,P\_j\} = 1 for every k in U, and R = U union \{j\}, and a supplied current reset state v\_U:M\_\{|U$|\}-\rangle$A\_U is an extended isomorphism satisfying epsilon\_U, d\_U <= t <= e\_s2 and d\_U <= c\_0\textasciicircum{}cb*epsilon\_U, then lem-compcb-corner-algebra makes A\_R an extended epsilon\_\{A\_R\}-C*-algebra, lem-maincb-nested-corner-comparison makes P\_U\textasciicircum{}R, P\_j\textasciicircum{}R quantitative projections in A\_R, and together with the lem-maincb-outer-compression-transfer outer-compressed isomorphism they satisfy every def-extcb-datum clause - approximate complementarity to I\_\{A\_R\}, one-dimensional S\_\{P\_j\textasciicircum{}R\}, nonzero S\_\{P\_U\textasciicircum{}R,P\_j\textasciicircum{}R\}, and total error e = delta + epsilon\_\{A\_R\} - with e <= C\_s2*t, forming the explicit Stage-2 raw-call closed EXT-CB datum in A\_R. \steptag{validated} \\[6pt]
\step{1.1} Universal constant ledger (excluding the unquantified additivity input). With C\_nest,e\_nest, e\_ncd, C\_out,e\_out, e\_sim, and C\_ca,e\_ca denoting the universal witnesses in the correspondingly named registered externals, set C\_delta=max\{1,C\_nest,C\_out\}, C\_s2=max\{1,C\_delta+2*C\_ca\}, and e\_s2=min\{1/2,e\_nest,e\_ncd,e\_out,e\_sim,e\_ca/2\}. These are positive finite universal constants, independent of m,U,j and every amplification. If t<=e\_s2 and epsilon\_A<=t, then t+epsilon\_A<=2*t<=e\_ca, and the named nested-comparison, nested-dimension-transport, outer-transfer, corner-equivalence, and corner-algebra thresholds are met; in particular the e\_ca threshold and C\_ca coefficient are absorbed exactly as the root requires. No applicability threshold for lem-extcb-corner-dimension-additivity is asserted or absorbed here, because its registered contract supplies no dimension-independent quantitative threshold. \steptag{validated} \\[6pt]
\step{1.1.1} Justification and scope of the corrected ledger. The registered externals explicitly provide positive universal e\_nest,e\_ncd,e\_out,e\_sim,e\_ca and finite universal C\_nest,C\_out,C\_ca. A finite minimum of the displayed positive thresholds is positive and universal, and the displayed finite maxima are finite and universal. Thus t<=e\_s2 implies each listed quantitative hypothesis; epsilon\_A<=t and e\_s2<=e\_ca/2 give t+epsilon\_A<=2*t<=e\_ca. In contrast, the registered contract of lem-extcb-corner-dimension-additivity says only sufficiently accurate and supplies neither a numeric threshold nor uniformity in m or |U|. Consequently it cannot enter this universal minimum, and this node gives no permission to invoke it downstream. The nonzero cross-corner clause must therefore be proved by a separate dimension-uniform argument or a strengthened allowed external; it does not follow from this ledger. \steptag{validated} \\[4pt]
\step{1.2} Original-corner geometry. Put e\_U=sum\_\{k in U\}e\_k, e\_R=e\_U+e\_j, P\_U=w(e\_U), and P\_R=w(e\_R)=P\_U+P\_j. Then P\_U,P\_j,P\_R are t-projections in A, P\_R is nonvanishing, every left/right subordination error of P\_U and P\_j to P\_R is at most t, A\_R=S\textasciicircum{}A\_\{P\_R\}, and the local definitions P\_U\textasciicircum{}R=Co\textasciicircum{}A\_\{P\_R\}(P\_U), P\_j\textasciicircum{}R=Co\textasciicircum{}A\_\{P\_R\}(P\_j) are the compressed elements required by the nested-corner externals. \steptag{validated} \\[6pt]
\step{1.2.1} Projection and nonvanishing check. At level one, def-extended-delta-inclusion makes w a t-homomorphism preserving involution and satisfying ||w(xy)-w(x)w(y)||<=t||x||||y|| together with (1-t)||x||<=||w(x)||<=(1+t)||x||. Since U and R=U union \{j\} are nonempty, e\_U,e\_j,e\_R are norm-one Hermitian projections of C\textasciicircum{}m. Hence their images P\_U,P\_j,P\_R are Hermitian and have idempotence defect at most t, so they are t-projections by def-delta-projection. Moreover 1-t<=||P\_R||<=1+t, whence | ||P\_R||-1 |<=t<=t+epsilon\_A; this is the nonvanishing alternative of def-delta-projection (with universal coefficient 1). \steptag{validated} \\[4pt]
\step{1.2.2} Subordination and corner notation. The coordinate identities e\_Ue\_R=e\_Re\_U=e\_U and e\_je\_R=e\_Re\_j=e\_j, inserted in the t-multiplicative-defect inequality, give ||P\_UP\_R-P\_U||, ||P\_RP\_U-P\_U||, ||P\_jP\_R-P\_j||, ||P\_RP\_j-P\_j||<=t. Linearity gives P\_R=P\_U+P\_j. The supplied def-maincb-partition-state defines A\_R=S\textasciicircum{}A\_\{P\_R\}; set P\_U\textasciicircum{}R=Co\textasciicircum{}A\_\{P\_R\}(P\_U) and P\_j\textasciicircum{}R=Co\textasciicircum{}A\_\{P\_R\}(P\_j), which are precisely the equalities demanded by lem-maincb-nested-corner-comparison and lem-maincb-nested-corner-dimension-transport. \steptag{validated} \\[4pt]
\step{1.3} Target algebra and projection clauses. Apply lem-compcb-corner-algebra to the nonvanishing t-projection P\_R in the extended epsilon\_A-C*-algebra A. It makes A\_R=S\textasciicircum{}A\_\{P\_R\}, with compressed product, inherited involution, and unit I\_\{A\_R\}=Co\textasciicircum{}A\_\{P\_R\}(P\_R), an extended epsilon\_\{A\_R\}-C*-algebra for epsilon\_\{A\_R\}=C\_ca*(t+epsilon\_A)<=2*C\_ca*t. Apply lem-maincb-nested-corner-comparison to (R,P,Q)=(P\_R,P\_U,P\_j), using node 1.2: P\_U\textasciicircum{}R and P\_j\textasciicircum{}R are C\_nest*t-projections in A\_R. \steptag{validated} \\[6pt]
\step{1.3.1} Corner-algebra application. Since epsilon\_A<=t, every defining extended epsilon\_A-C*-inequality also holds with the larger parameter t. Node 1.2 makes P\_R a nonvanishing t-projection. Node 1.1 gives t+epsilon\_A<=e\_ca. Hence lem-compcb-corner-algebra applies with delta=t and epsilon=epsilon\_A and gives A\_R=S\textasciicircum{}A\_\{P\_R\} the compressed product, inherited involution, and compressed unit u\_R=Co\textasciicircum{}A\_\{P\_R\}(P\_R), as an extended C\_ca*(t+epsilon\_A)-C*-algebra. Define epsilon\_\{A\_R\}=C\_ca*(t+epsilon\_A); then epsilon\_\{A\_R\}<=2*C\_ca*t and u\_R=I\_\{A\_R\}. \steptag{validated} \\[6pt]
\step{1.3.1.1} Local corner threshold, independent of node 1.1. The root conclusion is existential in the single universal witness e\_s2. Since the registered external lem-compcb-corner-algebra supplies a universal e\_ca>0, impose in constructing that witness the additional restriction 0<e\_s2<=e\_ca/2 (intersecting any other positive universal restrictions with e\_ca/2 preserves positivity). Then the root hypotheses epsilon\_A<=t<=e\_s2 give t+epsilon\_A<=2*t<=2*e\_s2<=e\_ca. Thus the external's required defect threshold holds without using node 1.1 or any additivity threshold. \steptag{validated} \\[4pt]
\step{1.3.1.2} Validated projection premise, independent of pending node 1.2. By validated node 1.2.1, P\_R is Hermitian, ||P\_R\textasciicircum{}2-P\_R||<=t, and 1-t<=||P\_R||<=1+t. Hence P\_R is a t-projection and satisfies the nonvanishing alternative | ||P\_R||-1 |<=t<=t+epsilon\_A. This is exactly the nonvanishing t-projection premise required by lem-compcb-corner-algebra. \steptag{validated} \\[4pt]
\step{1.3.1.3} Corner-algebra conclusion from the now explicit premises. The ambient A is finite-dimensional and is an extended epsilon\_A-C*-algebra by the root hypotheses. Child 1.3.1.1 proves t+epsilon\_A<=e\_ca, and child 1.3.1.2 proves that P\_R is a nonvanishing t-projection. Apply the registered lem-compcb-corner-algebra with delta=t, epsilon=epsilon\_A, and P=P\_R. It equips A\_R=S\textasciicircum{}A\_\{P\_R\} with the compressed product, inherited involution, and compressed unit u\_R=Co\textasciicircum{}A\_\{P\_R\}(P\_R), and makes it an extended C\_ca*(t+epsilon\_A)-C*-algebra. Set epsilon\_\{A\_R\}=C\_ca*(t+epsilon\_A) and I\_\{A\_R\}=u\_R. Since epsilon\_A<=t, epsilon\_\{A\_R\}<=2*C\_ca*t. Thus every conclusion in node 1.3.1 follows using children 1.3.1.1 and 1.3.1.2, with no reliance on pending nodes 1.1 or 1.2. \steptag{validated} \\[4pt]
\step{1.3.2} Nested-projection application. Regard A as an extended t-C*-algebra by the same parameter monotonicity. Node 1.2 supplies t-projections P\_R,P\_U,P\_j, nonvanishing P\_R, all four subordination errors at most t, A\_R=S\textasciicircum{}A\_\{P\_R\}, and the exact compressed-element definitions. Since node 1.1 gives t<=e\_nest, lem-maincb-nested-corner-comparison applies to (R,P,Q)=(P\_R,P\_U,P\_j) and concludes that P\_U\textasciicircum{}R and P\_j\textasciicircum{}R are C\_nest*t-projections in A\_R. Together with the first child this proves node 1.3. \steptag{validated} \\[6pt]
\step{1.3.2.1} Direct nested-comparison threshold bridge, with the final witness distinguished from the preliminary ledger. Let e\_s2\textasciicircum{}0 := min\{1/2, e\_nest, e\_ncd, e\_out, e\_sim, e\_ca/2\}. The fixed-source additivity repair in validated node 1.5.2 supplies a further positive universal threshold e\_\{2,1\} and chooses the final root witness e\_s2 := min\{e\_s2\textasciicircum{}0,e\_\{2,1\}\}. Consequently e\_s2 <= e\_s2\textasciicircum{}0 by the defining property of a minimum, and e\_s2\textasciicircum{}0 <= e\_nest because e\_nest is one of the entries defining e\_s2\textasciicircum{}0. Thus the root hypothesis t <= e\_s2 gives t <= e\_s2 <= e\_s2\textasciicircum{}0 <= e\_nest, so the nested-comparison threshold is met. The additivity threshold is used only to shrink the final universal witness; this nested-comparison application itself uses no conclusion of lem-extcb-corner-dimension-additivity. \steptag{validated} \\[6pt]
\step{1.3.2.1.1} Threshold derivation. Validated node 1.5.2 fixes the final root witness as e\_s2=min\{e\_s2\textasciicircum{}0,e\_\{2,1\}\}, where e\_s2\textasciicircum{}0 is the six-term minimum from node 1.1 and e\_\{2,1\}>0 is the fixed-(C\textasciicircum{}2,C) additivity threshold. Hence e\_s2<=e\_s2\textasciicircum{}0. From the definition e\_s2\textasciicircum{}0=min\{1/2,e\_nest,e\_ncd,e\_out,e\_sim,e\_ca/2\}, one also has e\_s2\textasciicircum{}0<=e\_nest. Combining these inequalities with the root assumption t<=e\_s2 yields t<=e\_nest. This proves exactly the threshold premise needed for lem-maincb-nested-corner-comparison without asserting the false equality e\_s2=e\_s2\textasciicircum{}0. \steptag{validated} \\[6pt]
\step{1.3.2.1.1.1} Dependency-explicit threshold chain. Validated node 1.1 defines the preliminary ledger witness e\_s2\textasciicircum{}0=min\{1/2,e\_nest,e\_ncd,e\_out,e\_sim,e\_ca/2\}; because e\_nest is an entry of this finite minimum, e\_s2\textasciicircum{}0<=e\_nest. Validated node 1.5.2 defines the final witness e\_s2=min\{e\_s2\textasciicircum{}0,e\_\{2,1\}\}, so e\_s2<=e\_s2\textasciicircum{}0. Therefore, under the root hypothesis t<=e\_s2, transitivity gives t<=e\_s2<=e\_s2\textasciicircum{}0<=e\_nest, exactly the threshold required by lem-maincb-nested-corner-comparison. \steptag{validated} \\[4pt]
\step{1.4} Diagonal isomorphism clause. The current reset state supplies v\_U:M\_|U$|-\rangle$A\_U=S\textasciicircum{}A\_\{P\_U\} as an extended d\_U-isomorphism. Since d\_U<=t, monotonicity of the defining defect and norm bounds makes it an extended t-isomorphism. Apply lem-maincb-outer-compression-transfer with R=P\_R and P=P\_U: the explicit map T=Co\textasciicircum{}\{A\_R\}\_\{P\_U\textasciicircum{}R\} o Co\textasciicircum{}A\_\{P\_R\} o v\_U is an extended C\_out*t-isomorphism M\_|U$|-\rangle$S\textasciicircum{}\{A\_R\}\_\{P\_U\textasciicircum{}R\}, and T\_n=I\_n tensor T for every n. The extra invariant d\_U<=c\_0\textasciicircum{}cb*epsilon\_U, with c\_0\textasciicircum{}cb supplied by lem-maincb-error-improvement, is retained as part of the input reset record but no improvement is needed in this conditional row. \steptag{validated} \\[4pt]
\step{1.5} Dimension clauses. One has dim S\textasciicircum{}\{A\_R\}\_\{P\_j\textasciicircum{}R\}=1 and S\textasciicircum{}\{A\_R\}\_\{P\_U\textasciicircum{}R,P\_j\textasciicircum{}R\} is nonzero. \steptag{validated} \\[6pt]
\step{1.5.1} Scalar corner. The phrase one-dimensional images P\_l means, by def-one-dimensional-delta-projection, that dim S\textasciicircum{}A\_\{P\_l\}=1; in particular dim S\textasciicircum{}A\_\{P\_j,P\_j\}=1. Apply lem-maincb-nested-corner-dimension-transport to (R,P,Q)=(P\_R,P\_j,P\_j). Node 1.2 supplies the t-projections, nonvanishing, both repeated left/right subordination bounds, A\_R=S\textasciicircum{}A\_\{P\_R\}, and P\_j\textasciicircum{}R=Co\textasciicircum{}A\_\{P\_R\}(P\_j), while node 1.1 supplies t<=e\_ncd. Therefore dim S\textasciicircum{}\{A\_R\}\_\{P\_j\textasciicircum{}R,P\_j\textasciicircum{}R\}=dim S\textasciicircum{}A\_\{P\_j,P\_j\}=1, i.e. dim S\textasciicircum{}\{A\_R\}\_\{P\_j\textasciicircum{}R\}=1. \steptag{validated} \\[4pt]
\step{1.5.2} Cross corner, using only a fixed-source additivity threshold. Choose k\_0 in the nonempty set U. If U=\{k\_0\}, then P\_U=P\_\{k\_0\}, so the hypothesis dim S\textasciicircum{}A\_\{P\_\{k\_0\},P\_j\}=1 already gives S\textasciicircum{}A\_\{P\_U,P\_j\} nonzero. Suppose instead that V=U\textbackslash\{\}\{k\_0\} is nonempty. Define exact non-unital *-monomorphisms alpha:C\textasciicircum{}2->C\textasciicircum{}m and beta:C->C\textasciicircum{}m by alpha(a,b)=a e\_\{k\_0\}+b e\_V and beta(c)=c e\_j. Because both coordinate blocks \{k\_0\},V are nonempty and disjoint, alpha is completely isometric; beta is completely isometric as well. Hence w o alpha and w o beta inherit, at every amplification, the t-homomorphism and two-sided (1+-t) norm bounds from the supplied non-unital extended t-inclusion w. Their units map respectively to P\_\{k\_0\}+P\_V=P\_U and P\_j, and their projection bases map to (P\_\{k\_0\},P\_V) and (P\_j). Apply lem-extcb-corner-dimension-additivity only to these fixed source algebras C\textasciicircum{}2 and C. Let e\_\{2,1\}>0 denote its sufficient-accuracy threshold for this fixed pair; because the pair is fixed, e\_\{2,1\} is a universal numerical constant independent of m,U,j and A. Write e\_s2\textasciicircum{}0 for the positive ledger constant constructed in node 1.1 and take the final root witness to be e\_s2=min\{e\_s2\textasciicircum{}0,e\_\{2,1\}\}. Every earlier deduction, which required only t<=e\_s2\textasciicircum{}0, remains valid under this shrink, while t<=e\_s2 now makes this fixed-source additivity call admissible. The resulting linear bijection is S\textasciicircum{}A\_\{P\_U,P\_j\} -> S\textasciicircum{}A\_\{P\_\{k\_0\},P\_j\} direct\_sum S\textasciicircum{}A\_\{P\_V,P\_j\}. Its first summand has dimension 1 by hypothesis, so S\textasciicircum{}A\_\{P\_U,P\_j\} is nonzero. In either case, lem-maincb-nested-corner-dimension-transport applied to (R,P,Q)=(P\_R,P\_U,P\_j), with the t-projections, nonvanishing and four subordination bounds from node 1.2 and t<=e\_ncd, gives dim S\textasciicircum{}\{A\_R\}\_\{P\_U\textasciicircum{}R,P\_j\textasciicircum{}R\}=dim S\textasciicircum{}A\_\{P\_U,P\_j\}>0. Thus the target cross corner is nonzero; no threshold uniform over variable-dimensional C\textasciicircum{}U and no equality dim S\textasciicircum{}A\_\{P\_U,P\_j\}=|U| is used. \steptag{validated} \\[4pt]
\step{1.6} Complementarity reduction and identified missing input. Validated node 1.2 gives P\_U+P\_j=P\_R and P\_U\textasciicircum{}R=Co\textasciicircum{}A\_\{P\_R\}(P\_U), P\_j\textasciicircum{}R=Co\textasciicircum{}A\_\{P\_R\}(P\_j); node 1.3 gives I\_\{A\_R\}=Co\textasciicircum{}A\_\{P\_R\}(P\_R). Therefore the desired equality P\_U\textasciicircum{}R+P\_j\textasciicircum{}R=I\_\{A\_R\} follows if Co\textasciicircum{}A\_\{P\_R\} is additive. But neither any definition registered in this workspace nor any allowed external asserts that additivity (or the defining compression formula from which it follows). Hence the former exact-complementarity claim is not derivable from the exact allowed inputs. A definition request has been filed for the compressed-corner compression map, including its defining formula and linearity/additivity; until that permitted input is provisioned, node 1.6 remains an explicit proof blocker and must not be used to discharge the complementarity clause. \steptag{validated} \\[4pt]
\step{1.7} Final assembly is blocked at complementarity. Let delta=C\_delta*t. If nodes 1.3-1.5 are validated, they provide the target extended epsilon\_\{A\_R\}-C*-algebra, the delta-projection clauses, the extended delta-isomorphism T:M\_|U$|-\rangle$S\textasciicircum{}\{A\_R\}\_\{P\_U\textasciicircum{}R\} with fixed amplifications, dim S\textasciicircum{}\{A\_R\}\_\{P\_j\textasciicircum{}R\}=1, and the nonzero cross-corner; the same constant calculation gives delta+epsilon\_\{A\_R\}<=(C\_delta+2*C\_ca)t<=C\_s2*t. However, node 1.6 does not prove ||P\_U\textasciicircum{}R+P\_j\textasciicircum{}R-I\_\{A\_R\}||<=delta: it proves only that exact complementarity would follow from additivity of Co\textasciicircum{}A\_\{P\_R\}, while also establishing that this additivity (or a substitute quantitative complementarity estimate) is absent from every registered definition and allowed external. Hence the exact allowed inputs do not establish the complementarity clause of def-extcb-datum, so they do not yet form a closed EXT-CB datum or a closed Stage-2 raw call. Closure requires provisioning an allowed definition/validated lemma giving compression additivity, or another allowed estimate implying ||P\_U\textasciicircum{}R+P\_j\textasciicircum{}R-I\_\{A\_R\}||<=delta; no such conclusion is claimed here. \steptag{validated} \\[4pt]
\step{1.8} Vacuity bridge from the permitted corner-equivalence input. Choose the root witness e\_s2 no larger than e\_sim (and, independently, no larger than every other positive universal threshold needed for the advertised absorption clauses; choose C\_s2 at least 1 and at least the finitely many advertised universal coefficients). Under the root hypotheses, epsilon\_A<=t makes A an extended t-C*-algebra by monotonicity of all defining inequalities, t<=e\_s2<=e\_sim, and every P\_l is a one-dimensional t-projection. Hence lem-maincb-corner-equivalence says that l\textasciitilde{}l' iff dim S\textasciicircum{}A\_\{P\_l,P\_l'\}=1 is an equivalence relation. By def-maincb-partition-state, the supplied current U is then a union of equivalence classes. Since U is nonempty, choose k in U. The hypothesis dim S\textasciicircum{}A\_\{P\_k,P\_j\}=1 says k\textasciitilde{}j, so j lies in the equivalence class of k. Because U is a union of equivalence classes and contains k, it contains that entire class, hence j in U, contradicting the root hypothesis j notin U. Thus no datum satisfies all root antecedents. \steptag{validated} \\[4pt]
\step{1.9} Final discharge by contradiction. Node 1.8 derives a contradiction solely from the root antecedents, def-maincb-partition-state, def-one-dimensional-delta-projection, and the allowed validated external lem-maincb-corner-equivalence, after shrinking the existential universal e\_s2 so that e\_s2<=e\_sim while retaining the advertised e\_ca/constant absorptions. Therefore the implication asserted by node 1 holds vacuously for those universal witnesses. In particular no complementarity estimate or additivity of Co\textasciicircum{}A\_\{P\_R\} is required; nodes 1.6 and 1.7 correctly identify a gap in the nonvacuous construction route but do not obstruct this contradiction proof of the root contract. \steptag{validated} \\[4pt]
\end{proofbox}

\end{document}
